\section{Conclusion}

In this paper, we propose Pre-hoc Sparsity (PrHS) framework, and derive an upper bound on the mutual-information (MI) loss based on the dropped attention mass. By tightening the MI bound towards the top-$k$ oracle, selectors under PrHS can achieve near-oracle KV compression. We propose \textbf{CPE} that instantiate three orthogonal modules under PrHS: \emph{Clustered Indices Sharing} (CIS), \emph{Progressive Sliding Attention Window} (PSAW), and \emph{Early Token Freezing} (ETF). CIS expands the highest-score indices in each shared critical set to its neighbors, achieving high recall and preserving accuracy even at aggressive (${\ge}90\%$) sharing ratios. PSAW and ETF further amplify sparsity by \textbf{masking tokens} that lie outside a dynamically sliding context window and \textbf{freezing tokens} whose updates have converged, thereby eliminating redundant computation with negligible performance loss. %These selectors are parameterized and tighten. 
Extensive evaluations on GSM8K, COQA, and LongBench demonstrate that CPE (i) matches dense baselines in accuracy under high attention sparsity, (ii) reduces retrieval complexity by up to $3\times$ than recent SOTA HShare~\cite{wuhshare} with even better performance, (iii) delivers up to a $9.9\times$ latency speed-up over FlashAttention and a $2.8\times$ throughput gain over GPT-Fast. These results confirm CPE as an effective and broadly applicable solution for efficient long-context inference in LLMs.


% In this paper, we propose HyperKV, which incorporates three key techniques: Clustered Indices Sharing (CIS), Progressive Sliding Attention Window (PSAW), and Early Token Freezing (ETF). CIS improves the retrieval rate of critical tokens by expanding shared critical indices to their neighbors, thereby maintaining strong model accuracy under high sharing ratios. PSAW and ETF enhance attention sparsity with minimal performance loss by eliminating redundant attention computations. We provide formal theoretical analyses to demonstrate the feasibility of these techniques and support their effectiveness through extensive experiments, showing that HyperKV consistently outperforms existing baselines.